\section*{Chapitre \ref{ch:b2:wittig} --- La réaction de Wittig et les autres formations de C=C}

\begin{solution}{exo:b2:wittig:1}
\ce{Ph3P + CH3CH2Br -> [Ph3PCH2CH3]+ + Br-} (substitution bimoléculaire, par chauffage dans un solvant),
puis butyllithium dans l'éther anhydre ou le THF sous gaz inerte~: \ce{[Ph3PCH2CH3]+ + BuLi -> Ph3P=CHCH3 + BuH + Li+}.
\end{solution}

\begin{solution}{exo:b2:wittig:2}
L'ylure \ce{Ph3P=CH2} et le benzaldéhyde donnent le phényléthène (styrène), \ce{C6H5CH=CH2}, et
l'oxyde de triphénylphosphine.
\end{solution}

\begin{solution}{exo:b2:wittig:3}
\ce{Ph3P=CHCOOEt} et \ce{Ph3P=CHCOCH3}~: stabilisés. \ce{Ph3P=CHCH3}~: non stabilisé, celui qui exige
le butyllithium (ou l'hydrure de sodium, l'amidure de sodium). \ce{Ph3P=CHC6H5}~: stabilisé par un aryle,
préparé en pratique avec un alcoolate ou un hydroxyde, et donnant des mélanges de géométries.
\end{solution}

\begin{solution}{exo:b2:wittig:4}
Le (\textit{E})-pent-2-énoate d'éthyle, \ce{CH3CH2CH=CHCOOEt}, avec du phosphate de diéthyle et de sodium.
\end{solution}

\begin{solution}{exo:b2:wittig:5}
Wittig~: l'ylure \ce{Ph3P=CH2}, préparé à partir de l'iodure de méthyltriphénylphosphonium et du
butyllithium, donne avec la cyclohexanone uniquement le méthylènecyclohexane, la double liaison étant
hors du cycle. La déshydratation du 1-méthylcyclohexan-1-ol passe par un carbocation tertiaire et perd
l'hydrogène qui donne l'alcène le plus substitué (Zaïtsev)~: surtout le 1-méthylcyclohexène, le mauvais isomère.
\end{solution}

\begin{solution}{exo:b2:wittig:6}
Les deux coupures sont identiques~: le benzaldéhyde et l'ylure de benzylidène \ce{Ph3P=CHC6H5}. Cet ylure
n'est stabilisé que par un aryle et donne un mélange (\textit{E})/(\textit{Z}). Pour l'isomère (\textit{E})
pur, utiliser la réaction HWE du benzylphosphonate de diéthyle avec le benzaldéhyde.
\end{solution}

\begin{solution}{exo:b2:wittig:7}
Avec \ce{Ph3P=CHCH3}~: le (\textit{Z})-pent-2-ène, \ce{CH3CH2CH=CHCH3}, surtout (\textit{Z}). Avec
\ce{Ph3P=CHCOOEt}~: le (\textit{E})-pent-2-énoate d'éthyle, surtout (\textit{E}).
\end{solution}

\begin{solution}{exo:b2:wittig:8}
\ce{Ph3P=CH2 + C6H10O -> C7H12 + Ph3PO}. Masses molaires~: méthylènecyclohexane \qty{96.173}{g/mol},
oxyde de triphénylphosphine \qty{278.291}{g/mol}~; économie d'atomes $96.173/(96.173 + 278.291) =
96.173/374.464 \approx 25.7~\%$. Les trois quarts de la masse des réactifs partent sous forme d'oxyde de phosphine.
% ledger: aw:C, aw:H, aw:O, aw:P
\end{solution}

\begin{solution}{exo:b2:wittig:9}
Le phosphore, nucléophile, déplace le bromure~:
\begin{equation*}
\ce{P(OEt)3 + BrCH2COOEt -> [(EtO)3PCH2COOEt]+ + Br-}.
\end{equation*}
Le bromure
attaque ensuite un carbone d'un groupe éthyle de l'ion phosphonium (substitution bimoléculaire), ce qui
libère le bromoéthane et forme la liaison \ce{P=O}~:
\begin{equation*}
\ce{[(EtO)3PCH2COOEt]+ + Br- -> (EtO)2P(O)CH2COOEt + EtBr}.
\end{equation*}
\end{solution}

\begin{solution}{exo:b2:wittig:10}
\ce{OHC-CH=CH-CHO + 2Ph3P=CHC6H5 -> C6H5CH=CH-CH=CH-CH=CHC6H5 + 2Ph3PO}. Il faut deux équivalents de
l'ylure de benzylidène (issu du chlorure de benzyltriphénylphosphonium et d'une base), un pour chaque
groupe aldéhyde. Les nouvelles doubles liaisons se forment en mélanges de géométries~; l'isomère
tout-(\textit{E}), le plus stable, est isolé par cristallisation, ou bien le mélange est isomérisé.
\end{solution}

\begin{solution}{exo:b2:wittig:11}
\omterm{def:b2:enolates-aldol:aldol}{Crotonisation}~: le benzaldéhyde avec un excès de propanone et l'hydroxyde de sodium (\cref{ch:b2:enolates-aldol})~;
réactifs bon marché, eau comme sous-produit, produit (\textit{E}), mais la condensation peut se produire
deux fois (dibenzalacétone) si la propanone n'est pas en excès. Wittig~: le benzaldéhyde avec l'\omterm{def:b2:wittig:ylide}{ylure
stabilisé} \ce{Ph3P=CHCOCH3}~; (\textit{E}), un seul produit, pas de double condensation, mais l'ylure
est coûteux et il faut éliminer l'oxyde de triphénylphosphine. La \omterm{def:b2:enolates-aldol:aldol}{crotonisation} est ici la meilleure voie.
\end{solution}

\begin{solution}{exo:b2:wittig:12}
Wittig~: le butanal avec \ce{Ph3P=CHCH2CH2CH3} (issu du 1-bromobutane, puis du butyllithium), en
l'absence de sels~: surtout (\textit{Z})~; sous-produits~: l'oxyde de triphénylphosphine et un peu
d'isomère (\textit{E}). Alcyne~: l'oct-4-yne (issu de l'éthyne par deux alkylations par le
1-bromopropane et l'amidure de sodium) hydrogéné sur un catalyseur au palladium empoisonné
(\cref{prop:b2:alkene-redox:syn-hydrogenation})~: uniquement (\textit{Z})~; les sous-produits sont le
bromure de sodium et l'ammoniac des alkylations.
\end{solution}

\begin{solution}{pb:b2:wittig:1}
\textbf{1.} \ce{C23H46}, soit \ce{C_nH_{2n}} avec $n = 23$~: une double liaison (ou un cycle), ici une
\ce{C=C}.
\textbf{2.} La double liaison \ce{C9=C10}.
\textbf{3.} Le nonanal \ce{CH3(CH2)7CHO} avec l'ylure issu d'un 1-halogénotétradécane~; ou le tétradécanal
\ce{CH3(CH2)12CHO} avec l'ylure issu d'un 1-halogénononane.
\textbf{4.} Le 1-bromotétradécane, \ce{CH3(CH2)13Br}.
\textbf{5.} La substitution bimoléculaire par la triphénylphosphine, volumineuse, est rapide sur un
carbone primaire non encombré et ne subit pas la concurrence de l'élimination.
\textbf{6.} \ce{Ph3P + C14H29Br -> [Ph3PC14H29]+ + Br-}, le bromure de tétradécyltriphénylphosphonium.
\textbf{7.} Le butyllithium (ou l'hydrure de sodium, l'amidure de sodium)~; l'hydroxyde est une base trop
faible pour un ylure non stabilisé, et l'eau qu'il apporte protonerait l'ylure.
\textbf{8.} Non, il ne porte qu'un groupe alkyle~: surtout (\textit{Z}) (\cref{prop:b2:wittig:stereo}).
\textbf{9.} Un cycle à quatre atomes \ce{P-C-C-O}~: le phosphore lié au carbone qui porte la chaîne
\ce{C13H27}, ce carbone lié à celui qui porte la chaîne \ce{C8H17}, et celui-ci lié à l'oxygène, lui-même
lié au phosphore.
\textbf{10.} Le butyllithium et l'ylure réagissent avec l'eau et le dioxygène.
\textbf{11.} La double liaison relie l'ancien carbone du carbonyle du nonanal (C9 du produit) à l'ancien
carbone de l'ylure (C10), et aucune autre liaison ne bouge (\cref{prop:b2:wittig:regiospecific}).
\textbf{12.} Un carbocation en C9, qui peut perdre un hydrogène de C8 ou de C10 et peut migrer~: un mélange
de tricos-8-ène et de tricos-9-ène, chacun (\textit{E}) et (\textit{Z}), surtout (\textit{E}).
\textbf{13.} \ce{Ph3P=CHC13H27 + C8H17CHO -> C8H17CH=CHC13H27 + Ph3PO}.
\textbf{14.} Cible \ce{C23H46}~: \qty{322.621}{g/mol}~; nonanal \ce{C9H18O}~: \qty{142.242}{g/mol}~; sel
\ce{C32H44BrP}~: \qty{539.582}{g/mol}~; \ce{Ph3PO}, \ce{C18H15OP}~: \qty{278.291}{g/mol}.
\textbf{15.} L'ylure \ce{C32H43P} a pour masse molaire $539.582 - 80.912 = \qty{458.670}{g/mol}$~; avec le
nonanal, \qty{600.912}{g/mol} de réactifs pour \qty{322.621}{g/mol} de produit~: $322.621/600.912 \approx
53.7~\%$.
\textbf{16.} Par sa polarité~: l'hydrocarbure, apolaire, traverse une courte colonne de silice avec un
éluant apolaire tandis que l'oxyde reste fixé~; ou bien l'oxyde cristallise dans un solvant apolaire froid.
\textbf{17.} La liaison \ce{P=O} formée, très forte, l'emporte sur la liaison \ce{P-C} rompue
(\cref{prop:b2:wittig:driving}).
\textbf{18.} Elle exige un \omterm{def:b2:wittig:hwe}{phosphonate} portant un groupe attracteur d'électrons sur le carbone qui réagit,
absent d'une simple chaîne alkyle, et elle donne l'alcène (\textit{E}).
\textbf{19.} Le tricos-9-yne, \ce{CH3(CH2)7C#C(CH2)12CH3}, avec le dihydrogène sur un catalyseur au
palladium empoisonné (palladium sur carbonate de calcium avec un sel de plomb et de la quinoléine).
\textbf{20.} Le déc-1-yne, \ce{CH3(CH2)7C#CH}, déprotoné par l'amidure de sodium, puis le 1-bromotridécane
\ce{CH3(CH2)12Br}~: $8 + 2 + 13 = 23$ carbones, la triple liaison entre C9 et C10.
\textbf{21.} Une hydrogénation plus poussée donnerait le tricosane, sans double liaison~; le catalyseur
empoisonné fixe l'alcyne bien plus fortement que l'alcène, qui quitte la surface avant une seconde addition.
\textbf{22.} Les deux voies comptent deux étapes à partir de fragments disponibles. La voie par l'alcyne
donne proprement l'isomère (\textit{Z}) et ne laisse que des sels~; la voie de Wittig donne un peu
d'isomère (\textit{E}) et une masse importante d'oxyde de triphénylphosphine.
\textbf{23.} 100~\%~: \ce{C23H44 + H2 -> C23H46}, chaque atome se retrouve dans le produit.
\textbf{24.} $278.291/322.621 = \boldsymbol{\approx \qty{0.86}{g}}$ d'oxyde de triphénylphosphine par
gramme de phéromone.
\end{solution}
